\title{Direct Preference Optimization: Your Language Model is Secretly a Reward Model}

\begin{abstract}
Large-scale unsupervised language models (LMs) acquire extensive world knowledge and some reasoning abilities during training, but fine-grained regulation of their output remains challenging, owing to the unsupervised structure of their learning process. To address this limitation, current approaches rely on human-annotated data that ranks model outputs, and then adjust the LM via fine-tuning to better match these rankings, often through reinforcement learning from human feedback (RLHF). RLHF, however, introduces complexity and instability by requiring the construction of a reward model mirroring human preferences, and subsequently fine-tuning the original LM through reinforcement learning in a way that maximizes the predicted reward while maintaining proximity to the base model. This work proposes a novel way to parameterize the reward model within RLHF, enabling direct retrieval of the optimal policy through a closed-form solution. This reformulation allows the classic RLHF setup to be solved simply using a classification loss, culminating in our proposed method: \textit{Direct Preference Optimization} (DPO). DPO delivers stable performance and strong results, while being computationally efficient—it bypasses the need for sampling from the LM during fine-tuning and obviates the necessity for extensive hyperparameter adjustments. Empirical evaluation demonstrates that DPO can align LMs with human preferences on par with, or superior to, existing techniques. In particular, DPO fine-tuning surpasses PPO-based RLHF in controlling sentiment of generated text, and either matches or surpasses it in summarization and single-turn dialogue tasks—all while requiring far less complexity in both implementation and training.
\end{abstract}

\section{Introduction}
Large unsupervised language models (LMs), when trained on extensive datasets, demonstrate unexpected emergent capabilities~\citep{chowdhery2022palm, brown2020language, touvron2023llama,bubeck2023sparks}. Nevertheless, the data underpinning these models is authored by humans whose objectives, levels of expertise, and values span a broad spectrum. Not all of these attributes are appropriate for LMs to reproduce: for instance, while it is useful for an AI coding assistant to \textit{recognize} typical programming errors in order to fix them, we generally wish for the model to favor the (sometimes infrequent) instances of superior coding present in its corpus when creating code. Analogously, though the language model ought to \textit{know} about misconceptions that are widespread (such as those held by 50\% of people), we certainly do not wish it to state such misconceptions as facts in half of the relevant outputs. In summary, precisely specifying the model’s \emph{target behaviors and responses}—given its extensive array of \textit{knowledge and skills}—is essential to achieving safe, reliable, and controllable AI systems \citep{ouyang2022training}. While common techniques guide LMs toward aligning with human preferences through reinforcement learning (RL), this work demonstrates that the objective underlying RL-based preference alignment can, in fact, be achieved directly with a straightforward binary cross-entropy loss, offering a substantial simplification to the standard preference learning workflow.

Generally, prior approaches aim to encode preferred behaviors in language models by utilizing carefully constructed datasets of human preferences, identifying traits that are safe and helpful from a human perspective. This process of preference learning typically takes place after the LM has undergone an initial phase of unsupervised pre-training on large-scale text. While a basic form of preference learning is to apply supervised fine-tuning using expert demonstrations, the most prevalent and effective strategies are based on reinforcement learning from human (or artificial) feedback (RLHF/RLAIF; \citep{christiano2017deep,bai2022constitutional}). RLHF procedures involve training a reward model to mimic human preferences and subsequently using RL to adjust the LM's outputs to maximize the reward signal while avoiding significant deviation from the pre-trained model. Although RLHF yields models with outstanding performance in dialogue and programming, it also introduces substantial complexity compared to straightforward supervised learning: the RLHF workflow necessitates the training of several models, and incorporates LM policy sampling within the training loop, resulting in notable computational overhead.

In this work, we introduce a method for optimizing language models to align with human preferences without relying on explicit reward models or reinforcement learning frameworks. We present \textit{{Direct Preference Optimization} (DPO)}, an algorithm that implicitly targets the same objective pursued by typical RLHF techniques—maximizing reward under a KL-divergence constraint—while maintaining simplicity in both its formulation and implementation. The core idea behind the DPO update is to enhance the log-likelihood ratio in favor of responses that are preferred by humans over those that are not. This approach employs a dynamically calculated importance weight for each example, which mitigates the risk of model collapse that can arise when using a naive probability ratio objective. DPO, like previous approaches, is grounded in a theoretical preference model (e.g., the Bradley-Terry model; \cite{bradley1952rankanalysis}) to evaluate how well a reward function captures real-world preference data. Yet, in contrast to standard methodologies—where the preference model is used to construct a loss for training a reward function, which then informs the training of the policy—DPO introduces a reparameterization, defining the preference loss directly in terms of the policy itself. As a result, given annotated human preference data over model outputs, DPO allows one to optimize the policy directly using a simple binary cross-entropy loss, achieving the optimal policy relative to an implicit reward shaped by the preference dataset.

The principal contribution of this paper is the proposal of Direct Preference Optimization (DPO), an RL-free and straightforward approach to training language models using preference data. Our empirical findings demonstrate that DPO matches or surpasses the performance of contemporary methods, such as PPO-based RLHF, in preference-driven learning tasks including sentiment adjustment, summarization, and conversational modeling, and does so for language models containing up to 6B parameters.

Self-supervised language models at ever-greater scales are able to tackle certain tasks without any task-specific data (zero-shot) \citep{radford2019language} or with minimal examples provided in the prompt (few-shot) \citep{gpt3,megatron,chowdhery2022palm}. Yet, their efficacy on downstream applications and their alignment with user intentions see substantial gains when they are further trained—fine-tuned—using collections of instructional prompts paired with human-generated responses \citep{mishra-etal-2022-cross,sanh2022multitask,chung2022scaling,thoppilan2022lamda}. This process, termed `instruction-tuning,' equips large language models (LLMs) to adapt to novel instructions not seen during tuning and enhances their overall practical value \citep{chung2022scaling}. While instruction tuning has achieved notable results, acquiring \textit{comparative} human evaluations of model responses is typically more feasible than sourcing high-quality expert data, prompting recent research to fine-tune LLMs based on datasets reflecting human preference judgments. This approach has led to improvements in tasks including translation \citep{kreutzer-etal-2018-reliability}, summarization \citep{stiennon2022learning,ziegler2020finetuning}, narrative generation \citep{ziegler2020finetuning}, and executing instructions \citep{ouyang2022training,ramamurthy2023is}. In these frameworks, a neural reward model is first trained to predict which response aligns better with the recorded preferences, using models like the Bradley-Terry preference model \citep{bradley1952rankanalysis}. Next, reinforcement learning algorithms—typically REINFORCE \citep{williams1992reinforce}, proximal policy optimization (PPO; \cite{schulman2017proximal}), or related techniques \citep{ramamurthy2023is}—are employed to optimize the language model with respect to this learned reward. Related advancements utilize instruction-tuned LLMs, further refined with human feedback, to create new synthetic preference data targeting properties such as safety or harmlessness \citep{bai2022constitutional}; here, the LLM uses only basic text-based guidelines from humans to label data. These efforts unite two principal research directions: the use of reinforcement learning for various language model training objectives~\citep{Ranzato2015SequenceLT,paulus2018a,wu2018learning}, and broader strategies for leveraging human preference information \citep{christiano2017deep,kupcsik2018learning}. Even though utilizing comparative human preference data has clear practical benefits, reinforcement learning-based fine-tuning of large language models poses significant practical obstacles; this study introduces a theoretically grounded method to optimize with relative preferences that circumvents reinforcement learning entirely.

Beyond linguistic applications, the problem of learning policies from preferences has been explored within both bandit and reinforcement learning domains, resulting in several distinct methodologies. The contextual dueling bandit (CDB; \cite{yue2012karmed,dudik2015contextual}) paradigm, for instance, substitutes conventional reward signals with ranked or preferred actions as supervisory signals. In theoretical analyses, due to the lack of explicit reward feedback, CDBs conceptualize an optimal policy as a \textit{von Neumann winner}: a policy whose expected head-to-head victory rate versus \textit{any} alternative policy meets or exceeds 50\% \citep{dudik2015contextual}. Notably, CDB settings assume preferences are provided interactively, while learning from human preference data is more often carried out over a static, pre-collected set of preference-labeled action pairs \citep{yan2022human}. In a similar vein, \textit{preference-based RL} (PbRL) leverages binary feedback arising from an unspecified underlying `scoring' function, rather than direct reward observations \citep{BusaFekete2014,ruiz2023dueling}. Numerous approaches have been developed within PbRL, encompassing techniques capable of exploiting off-policy preference datasets; most, however, explicitly reconstruct a latent scoring or reward function prior to policy optimization \citep{jain2013learning,BusaFekete2014,christiano2017deep,sadigh2017active,kupcsik2018learning}. Here, we propose an alternative: a unified stage of policy learning that seeks to directly align a policy with observed preferences.

\section{Preliminaries}\label{section:prelims}

We briefly summarize the RLHF framework as presented by \citeauthor{ziegler2020finetuning} and subsequently expanded upon in works such as \citep{stiennon2022learning, bai2022training, ouyang2022training}. This process commonly comprises three main steps: 1) supervised fine-tuning (SFT); 2) collection of preferences and construction of a reward model; and 3) reinforcement learning for optimization.

\textbf{SFT}: The initial phase of RLHF involves adapting a pre-trained language model through supervised training on a curated dataset relevant to the targeted task(s)—such as dialogue or summarization—yielding the fine-tuned model $\pi^\text{SFT}$.

\textbf{Reward Modelling Phase}: The subsequent stage uses the SFT model to generate response pairs in reaction to a set of prompts $x$, resulting in $(y_1, y_2)\sim \pi^\text{SFT}(y \mid x)$. Human annotators then assess these output pairs by indicating a preference—for example, expressing that $y_w\succ y_l \mid x$ where $y_w$ denotes the favored and $y_l$ the less favored completion in $(y_1, y_2)$. It is assumed these preferences reflect evaluations from some unknown reward mechanism $r^*(y, x)$ that is not directly observable. Several probabilistic models exist to capture this preference data, with the Bradley-Terry (BT) \cite{bradley1952rankanalysis} framework being especially prevalent (although more general models such as the Plackett-Luce ranking model \citep{plackett1975analysis, luce2012individual} can also be utilized, particularly if ranked lists rather than pairs are available). Under the BT model, the human-generated preference distribution $p^*$ is expressed as follows:

Given a static set of comparison data $\mathcal{D}=\bigl\{x^{(i)}, y_w^{(i)}, y_l^{(i)}\bigr\}_{i=1}^N$ sampled from $p^*$, one may model the reward function $r_{\phi}(x, y)$ and estimate its parameters through maximum likelihood estimation. Viewing this as a binary classification setup, the negative log-likelihood objective is given by:

\begin{equation}\label{eq:reward_model}
    \mathcal{L}_R(r_{\phi}, \mathcal{D}) = -\mathbb{E}_{(x, y_w, y_l)\sim \mathcal{D}}\left[\log \sigma(r_{\phi}(x, y_w) - r_{\phi}(x, y_l))\right]
\end{equation}

where $\sigma$ denotes the logistic sigmoid function. In the language modeling context, $r_{\phi}(x, y)$ is commonly initialized from the SFT model $\pi^\text{SFT}(y \mid x)$, augmented by appending a linear projection atop the transformer to yield a single scalar representing reward \cite{ziegler2020finetuning}. To regularize and control variance in the learned reward, it is typical to normalize such that $\mathbb{E}_{x,y\sim \mathcal{D}}[r_{\phi}(x, y)] = 0$ across the data distribution.

\textbf{RL Fine-Tuning Phase}: In the subsequent reinforcement learning stage, this reward estimator is used to provide signal for further updating the language model. In accordance with earlier works~\citep{jaques2017sequence, jaques2020human}, optimization proceeds as follows:

\begin{equation}\label{eq:RL}
\max_{\pi_{\theta}} \; \mathbb{E}_{x\sim \mathcal{D}, y\sim \pi_{\theta}(y \mid x)}\left[r_{\phi}(x, y)\right] - \beta\mathbb{D}_{\textrm{KL}}\left[\pi_{\theta}(y\mid x)\mid \mid \pi_\text{ref}(y\mid x)\right],
\end{equation}

with $\beta$ mediating the regularization towards a reference policy $\pi_\text{ref}$, typically chosen as the original SFT model $\pi^\text{SFT}$. The policy $\pi_\theta$ is generally initialized from this same SFT baseline. This regularization is crucial to constrain policy updates, ensuring generated samples remain within the distribution where the reward model remains reliable, as well as to sustain response diversity and preclude collapse to isolated modes. Owing to the inherent discreteness of natural language outputs, this objective cannot be directly differentiated and thus requires reinforcement learning methods. The prevailing strategy~\citep{ziegler2020finetuning, stiennon2022learning, bai2022training, ouyang2022training} is to employ the reward $r(x, y) = r_{\phi}(x, y) -\beta \left(\log \pi_{\theta}(y\mid x) - \log \pi_\text{ref}(y\mid x)\right)$ and optimize it via PPO~\cite{schulman2017proximal}. 

\section{Direct Preference Optimization}\label{sec:DPO}

Considering the computational and algorithmic hurdles encountered when applying RL algorithms to large language models, we seek a streamlined policy optimization method that acts directly on preference data. Departing from earlier RLHF paradigms, which first estimate a reward before performing RL-based optimization, our strategy exploits a tailored reward model structure that affords a closed-form solution for its corresponding optimal policy—bypassing the need for a separate RL procedure. The cornerstone of our method is an analytical relationship linking reward functions to their respective optimal policies, which makes it possible to transfer optimization from reward parameters to policy parameters directly. This transformation circumvents explicit training of an independent reward network, yet continues to optimize using existing models of human preference, such as the Bradley-Terry framework. Effectively, the policy model takes on the joint role of both language generator and implicit reward estimator.

\textbf{Derivation of the DPO objective.} We begin by considering the standard RL objective, as formulated in Eq.~\ref{eq:RL}, defined with respect to an arbitrary reward function $r$. Building on the results of earlier works~\citep{peters2007reinforcement, peng2019advantage, korbak2022reinforcement, go2023aligning}, it can be readily demonstrated that the optimal policy for the KL-regularized reward maximization problem in Eq.~\ref{eq:RL} is given by:

\begin{equation}\label{eq:op_policy}
    \pi_r(y\mid x) = \frac{1}{Z(x)}\pi_\text{ref}(y\mid x)\exp\left(\frac{1}{\beta}r(x, y)\right),
\end{equation}

where the partition function is $Z(x) =\sum_{y}\pi_\text{ref}(y\mid x)\exp\left(\frac{1}{\beta}r(x, y)\right)$. Full details of this derivation are provided in Appendix~\ref{app:derivation1}. Although employing an MLE estimator $r_{\phi}$ for the true reward $r^*$ is possible, in practice, computation of $Z(x)$ remains computationally intensive~\citep{korbak2022reinforcement, go2023aligning}, making this form of the solution less practical. Nevertheless, it is possible to rewrite Eq.~\ref{eq:op_policy} to isolate the reward as a function of the optimal policy $\pi_r$, the baseline policy $\pi_\text{ref}$, and the partition function $Z(\cdot)$. Specifically, by taking the natural logarithm of both sides of Eq.~\ref{eq:op_policy} and rearranging, we obtain:

\begin{equation}\label{eq:main_eq}
    r(x,y) =\beta \log \frac{\pi_r(y\mid x)}{\pi_\text{ref}(y\mid x)} + \beta \log Z(x).
\end{equation}

Applying this transformation to the true reward $r^*$ and its corresponding optimal policy $\pi^*$ is straightforward. Importantly, the Bradley-Terry framework depends solely on the reward difference between two candidate completions: ${p^*(y_1 \succ y_2 \mid x) = \sigma(r^*(x, y_1) - r^*(x, y_2))}$. If we substitute the expression for $r^*(x, y)$ from Eq.~\ref{eq:main_eq} into the preference probability defined in Eq.~\ref{eq:bradley-terry}, the partition terms $Z(x)$ cancel out, leaving the human preference probabilities dependent only on the optimal policy $\pi^*$ and reference policy $\pi_\text{ref}$. This establishes that under the Bradley-Terry model, the RLHF-optimal policy $\pi^*$ is characterized by:

\begin{equation}\label{eq:objective}
    p^*(y_1\succ y_2 \mid x)=\frac{1}{1 + \exp\left(\beta \log \frac{\pi^*(y_2\mid x)}{\pi_\text{ref}(y_2\mid x)} - \beta \log \frac{\pi^*(y_1\mid x)}{\pi_\text{ref}(y_1\mid x)}\right)}
\end{equation}

The detailed proof can be found in Appendix~\ref{app:derivation2}. While Eq.~\ref{eq:objective} is framed using the Bradley-Terry model, analogous results for the broader Plackett-Luce models~\citep{plackett1975analysis, luce2012individual} are presented in Appendix~\ref{app:plackett_luce_models}.

Having established a formulation of human preference probabilities in terms of the optimal policy as opposed to the reward function, we are positioned to formulate a maximum likelihood objective over a parametric policy $\pi_\theta$. Mirroring the approach in reward modeling (see Eq.~\ref{eq:reward_model}), we define our policy learning objective as:

\begin{equation}\label{eq:optimum_model}
    \mathcal{L}_\text{DPO}(\pi_{\theta}; \pi_\text{ref}) = -\mathbb{E}_{(x, y_w, y_l)\sim \mathcal{D}}\left[\log \sigma \left(\beta \log \frac{\pi_{\theta}(y_w\mid x)}{\pi_\text{ref}(y_w\mid x)}

In this framework, we model the implicit reward function through an alternative parameterization, such that the optimal policy naturally becomes $\pi_\theta$. Notably, our approach is tantamount to training a reparametrized Bradley-Terry model, granting it desirable theoretical guarantees like consistency, provided that the distribution of preference data satisfies certain conditions \cite{bong2022generalized}. Additional theoretical aspects of DPO, and its connections to related methodologies, are examined in Section~\ref{sec:theory}.

\textbf{What is the effect of the DPO update?} To gain a mechanistic perspective on DPO, it is instructive to consider the gradient of its objective $\mathcal{L}_\text{DPO}$. The gradient with respect to $\theta$ takes the following form:

\begin{multline*}\label{eq:gradient}
    \nabla_\theta \mathcal{L}_\text{DPO}(\pi_\theta;\pi_\text{ref}) = \\ -\beta\mathbb{E}_{(x, y_w, y_l) \sim \mathcal{D}} \left[\underbrace{\sigma(\hat{r}_\theta(x, y_l) - \hat{r}_\theta (x, y_w))}_\text{larger influence when the reward prediction is inaccurate}\left(\underbrace{\nabla_\theta\log \pi(y_w \mid x)}_\text{boost probability of $y_w$} - \underbrace{\nabla_\theta\log\pi(y_l \mid x)}_\text{reduce probability of $y_l$}\right)\right],
\end{multline*}

where $\hat{r}_\theta(x, y) = \beta \log \frac{\pi_\theta(y \mid x)}{\pi_\text{ref}(y \mid x)}$ represents the implicitly specified reward by the language model $\pi_\theta$ relative to the reference model $\pi_\text{ref}$ (see Section~\ref{sec:theory} for more details). Conceptually, the gradient of $\mathcal{L}_\text{DPO}$ works by encouraging the model to favor higher probabilities for preferred completions $y_w$, while discouraging less desirable completions $y_l$. Crucially, this effect is modulated according to the extent to which the implicit reward $\hat{r}_\theta$ incorrectly scores the less-preferred completion, as scaled by $\beta$—essentially, by how much the current reward misorders the outputs, thus reflecting the strength of the KL constraint. Our empirical findings highlight the necessity of this adaptive weighting; omitting the weighting factor, as in a straightforward version of the update, frequently results in model collapse (see Appendix Table~\ref{tab:unlikelihood_generations}).

\textbf{DPO overview.} The standard DPO workflow proceeds as follows: 1) For each prompt $x$, generate completions $y_1, y_2 \sim \pi_\text{ref}(\cdot \mid x)$, annotate them with human preference labels, and assemble the offline preference dataset $\mathcal{D} = \{x^{(i)}, y_w^{(i)}, y_l^{(i)}\}_{i=1}^N$; 2) Train the language model $\pi_\theta$ by minimizing $\mathcal{L}_\text{DPO}$ given $\pi_\text{ref}$, $\mathcal{D}$, and a selected $\beta$. 
In applied settings, researchers typically utilize pre-existing public preference datasets instead of producing new samples and collecting human annotations. Since these datasets are generally generated with $\pi^\text{SFT}$, it is common to set $\pi_\text{ref} = \pi^\text{SFT}$ when it is accessible. If $\pi^\text{SFT}$ is unavailable, $\pi_\text{ref}$ is instead initialized by maximizing the likelihood over preferred completions ${(x, y_w)}$ as follows: ${\pi_\text{ref} = \argmax_{\pi}\mathbb{E}_{x, y_w \sim \mathcal{D}}\left[\log \pi(y_w \mid x)\right]}$. This approach serves to counteract potential distributional discrepancies between the unknown true reference distribution and the practical $\pi_\text{ref}$ employed in DPO. Additional specifics on implementation choices and hyperparameter configurations are detailed in Appendix~\ref{app:implementation}.

\section{Theoretical Analysis of DPO}
This section delves into a deeper analysis of the DPO method, presenting theoretical justifications and highlighting the benefits of DPO in contrast to actor-critic approaches commonly adopted in RLHF, such as PPO~\cite{schulman2017proximal}.

\label{sec:theory}

\subsection{Your Language Model Is Secretly a Reward Model}
DPO circumvents both explicit reward modeling and policy learning through RL by leveraging a unified maximum likelihood objective. The training objective given in Eq. \ref{eq:main_eq} mirrors the Bradley-Terry framework when using the reward parametrization $r^*(x, y) = \beta \log\frac{\pi^*_\theta(y \mid x)}{\pi_\text{ref}(y \mid x)}$. Thus, optimizing the parameterized model $\pi_{\theta}$ becomes equivalent to the reward model training outlined in Eq. \ref{eq:reward_model}, once variables are appropriately substituted. In this section, we construct the theoretical underpinnings for this parameterization, demonstrate that it preserves the full representational capacity of reward models, and establish that the optimal policy can be exactly recovered. We begin by introducing an equivalence relation among reward functions.

\begin{definition}
Two reward functions $r(x, y)$ and $r'(x, y)$ are defined as equivalent if and only if there exists a function $f$ such that ${r(x, y)-r'(x, y) = f(x)}$.     
\end{definition}

One can readily verify that this definition indeed forms an equivalence relation, effectively dividing the set of reward functions into distinct equivalence classes. This allows us to articulate the following lemmas:

\begin{lemma}\label{lemma:same_prefrence}
Within the Plackett-Luce preference framework, including the Bradley-Terry model as a specific case, reward functions belonging to the same equivalence class generate identical preference distributions.
\end{lemma}

\begin{lemma}\label{lemma:same_policy}
    Any two reward functions within the same equivalence class yield the same optimal policy for the constrained RL problem.
\end{lemma}

As the proofs follow directly from standard arguments, we postpone them to Appendix \ref{app:lemma1}. The first lemma encapsulates the common under-specification property inherent to the Plackett-Luce family of models \cite{plackett1975analysis}. Because of this, it is necessary to enforce extra identifiability conditions in order to obtain guarantees about MLE solutions of Eq. \ref{eq:reward_model} \cite{bong2022generalized}. The second lemma establishes that all rewards from an equivalence class induce identical optimal policies, so for our purposes, it suffices to recover any reward representative from the optimal equivalence class. In Appendix~\ref{app:thm1}, we demonstrate the following theorem:

\begin{theorem}\label{thm:main}
    Subject to mild regularity conditions, every reward equivalence class consistent with the Plackett-Luce (and specifically the Bradley-Terry) models admits a representation of the form ${r(x, y) = \beta \log \frac{\pi(y\mid x)}{\pi_\text{ref}(y\mid x)}}$ for an appropriate model $\pi(y\mid x)$ and a chosen reference model $\pi_\text{ref}(y \mid x)$.
\end{theorem}

\begin{sproof}
    Take any reward function $r(x, y)$, which defines an associated optimal model $\pi_r(y \mid x)$ through Eq. \ref{eq:op_policy}. We aim to demonstrate that some member of the equivalence class of $r$ can be written using the above reparameterization. Define the mapping $f$ by  
\begin{equation}
    f(r; \pi_\text{ref}, \beta)(x, y) = r(x, y) - \beta\log\sum_{y}\pi_\text{ref}(y\mid x)\exp\left(\frac{1}{\beta}r(x, y)\right)
\end{equation}
This operator normalizes $r$ by subtracting the logarithm of the partition function with respect to $\pi_\text{ref}$. The normalization term depends only on $x$, so $f(r; \pi_\text{ref}, \beta)(x, y)$ remains in the equivalence class of $r(x, y)$. Substituting $r$ by the expression in Eq.~\ref{eq:main_eq} (which is valid for any reward), we observe that $f(r; \pi_\text{ref}, \beta)(x, y) = \beta \log \frac{\pi_r(y\mid x)}{\pi_\text{ref}(y\mid x)}$. Therefore, the function $f$ produces an equivalent reward function of the required form, and the reparameterization fully characterizes the set of reward models of interest.
\end{sproof}

An alternative perspective on Theorem~\ref{thm:main} is that it pinpoints, within each equivalence class, the specific reward function chosen by the DPO reparameterization. In particular, it identifies the reward function that satisfies

\begin{equation}\label{eq:lag_p}
     \sum_{y}\underbrace{\pi_\text{ref}(y\mid x)\exp\left(\frac{1}{\beta}r(x, y)\right)}_{=\pi(y\mid x)\text{, using Thm.~\ref{thm:main} reparam.}} = 1,
\end{equation}

which ensures that $\pi(y\mid x)$ defines a legitimate probability distribution: all probabilities are non-negative and the total sums to one. Moreover, by referring to Eq.~\ref{eq:op_policy}, we observe that Eq.~\ref{eq:lag_p} represents the partition function for the optimal policy specified by the reward function $r(x, y)$. The central idea underlying the DPO algorithm is to introduce particular constraints within the inherently underdetermined Plackett-Luce class (notably including the Bradley-Terry model) of preference models, thereby retaining the expressive capacity for reward functions but making the optimal policy in Eq.~\ref{eq:op_policy} computable in closed-form for any prompt $x$.

\subsection{Instability of Actor-Critic Algorithms}

Our framework additionally enables an analysis of instabilities encountered with prevalent actor-critic methods employed for RLHF, such as PPO. Adopting the RLHF protocol, our focus is on the RL fine-tuning stage as described in Section~\ref{section:prelims}. By relating this to the control as inference formulation \cite{levine2018reinforcement} in the context of the constrained RL problem from \ref{eq:RL}, we assume a parametric policy $\pi_{\theta}(y\mid x)$ and aim to minimize $\mathbb{D}_{\text{KL}}[\pi_{\theta}(y|x) \mid \mid \pi^*(y\mid x)]$, with $\pi^*$ denoting the optimal policy from Eq.~\ref{eq:optimum_model} defined via the reward function $r_{\phi}(y, x)$. After performing algebraic manipulation, the objective function can be written as:

\begin{equation}\label{eq:AC}
    \max_{\pi_{\theta}}\mathbb{E}_{\pi_{\theta}(y\mid x)}\bigg[\underbrace{r_{\phi}(x, y) -\beta\log\sum_{y}\pi_\text{ref}(y\mid x)\exp\left(\frac{1}{\beta}r_{\phi}(x, y)\right)}_{f(r_{\phi}, \pi_\text{ref}, \beta)} - \underbrace{\beta\log\frac{\pi_{\theta}(y\mid x)}{\pi_\text{ref}(y\mid x)}}_{\text{KL}}\bigg]
\end{equation}

The objective here mirrors that of earlier studies  
\citep{ziegler2020finetuning, stiennon2022learning, bai2022training, ouyang2022training}, employing the DPO-equivalent reward associated with the reward class $r_{\phi}$. Within this context, the normalization factor present in $f(r_{\phi}, \pi_\text{ref}, \beta)$ can be understood as the soft value function associated with the reference policy $\pi_\text{ref}$. Although this normalization does not alter the location of the optimum, omitting it can significantly increase the variance of the policy gradient, thereby reducing training stability. One possible solution involves approximating the normalization factor with a learned value function, but this approach introduces further optimization challenges. Alternatively, previous approaches often normalized rewards using a baseline derived from human completions, effectively constituting a one-sample Monte Carlo estimate of the normalization constant. By comparison, the DPO reparameterization offers a reward formulation that circumvents the necessity for any explicit baselines.

\section{Experiments}
Here, we provide empirical evidence regarding DPO’s capacity to learn policies solely from preference data. We start by investigating a controlled text generation environment to explore DPO’s effectiveness in balancing reward maximization and KL-divergence minimization with the reference policy. We directly compare its performance to established preference optimization methods such as PPO. Subsequently, we extend our analysis to include evaluations on larger model architectures and more complex RLHF challenges, for example, tasks involving summarization and conversational agents. Our findings reveal that, requiring minimal hyperparameter adjustment, DPO typically matches or surpasses the effectiveness of competitive methods like RLHF with PPO, and can outperform strategies based on selecting the best trajectory from $N$ candidates under a learned reward. Details of our experimental framework precede the presentation of these results, with further information available in Appendix~\ref{app:exp_details}.

\textbf{Tasks.} We investigate three distinct open-ended text generation tasks in our experiments. Across all settings, the methods learn a policy from a collection of preference-based samples denoted as $\mathcal{D}=\bigl\{x^{(i)}, y_w^{(i)}, y_l^{(i)}\bigr\}_{i=1}^N$. In the task of \textbf{controlled sentiment generation}, $x$ corresponds to a leading segment of a movie review sourced from the IMDb dataset \cite{maas-EtAl:2011:ACL-HLT2011}, and the objective for the policy is to generate continuations $y$ that express positive sentiment. To enable systematic assessment, we automatically create pairs of generations ranked by preference, utilizing an established sentiment classifier such that $p(\text{positive}\mid x,y_w)>p(\text{positive}\mid x,y_l)$. The SFT baseline is constructed by fine-tuning GPT-2-large until convergence on the IMDb dataset’s training reviews (refer to App~\ref{app:sentiment_details} for more information). In the \textbf{summarization} task, $x$ consists of Reddit forum posts, and the policy is tasked with composing a summary $y$ that encapsulates the main ideas of the post. We build on previous approaches, employing the Reddit TL;DR summarization dataset \citep{volske-etal-2017-tl} along with human-annotated preferences as collected by \citeauthor{stiennon2022learning}. For SFT in this context, a model pre-trained on summaries written by humans\footnote{\url{https://huggingface.co/CarperAI/openai_summarize_tldr_sft}} is utilized and adapted using the TRLX framework \citep{leandro_von_werra_2023_7790115} for RLHF. Human feedback used as preferences originates from the data gathered by \citeauthor{stiennon2022learning} on samples produced by another, similarly fine-tuned, SFT model. Lastly, for the \textbf{single-turn dialogue} task, $x$ represents a query from a user—ranging from scientific inquiries to requests for personal advice. The policy is required to generate a relevant, engaging, and useful response $y$; for this purpose, we leverage the Anthropic Helpful and Harmless dialogue dataset \citep{bai2022training}, which comprises 170k human–assistant conversation records. Each dialogue concludes with two model-generated responses, each associated with a human preference annotation to indicate the preferred reply. Due to the absence of an available pre-trained SFT model in this case, we instead fine-tune a general language model using only the completions labeled as preferred, thereby constructing the SFT baseline for this domain.

\textbf{Evaluation.} We adopt two primary strategies for evaluating our models. For the controlled sentiment generation experiments, where the ground-truth reward function is accessible (utilizing a sentiment classifier), we assess each algorithm by mapping the trade-off frontier between attained reward and KL-divergence relative to the reference policy. This is feasible given the known reward structure in this synthetic scenario. However, outside of controlled settings, true reward functions remain unknown. In these real-world tasks, we measure algorithm performance via their \textit{win rate} compared to a baseline policy. Specifically, we leverage GPT-4 as a stand-in for human judgment, evaluating both summary quality in summarization and helpfulness in single-turn dialogue contexts. For the summarization task, the baseline is set to the reference summary from the test corpus, while for dialogue, the preferred annotated response in the test set serves as the baseline. Previous literature indicates that LMs can surpass established metrics as automated evaluators \citep{Chen2023ExploringTU}, but we supplement this by carrying out a human study, validating our reliance on GPT-4 for evaluation as described in Sec.~\ref{sec:human-judgments}. Our findings reveal that GPT-4's ratings are highly aligned with human judgment, and agreement rates between humans and GPT-4 tend to match or exceed those found among human annotators themselves.

\textbf{Methods.} Alongside DPO, we benchmark a selection of prevailing techniques for tuning language models to match human preferences. For summarization, we employ zero-shot prompting using \textbf{GPT-J} \citep{gpt-j}, and for the dialogue task, we use 2-shot prompting with \textbf{Pythia-2.8B} \citep{biderman2023pythia}. Further, we test the \textbf{SFT} model and introduce \textbf{Preferred-FT}, a variant trained via supervised fine-tuning on the selected response $y_w$, sourced either from the SFT output (for sentiment control and summarization) or from a general language model (for single-turn dialogue). Another comparative approach is \textbf{Unlikelihood}~\citep{welleck2019neural}, which enforces the model to increase the likelihood of $y_w$ while simultaneously penalizing the likelihood assigned to $y_l$; here, an optional parameter $\alpha\in[0,1]$ can be applied to modulate the `unlikelihood' objective. Additionally, we include \textbf{PPO} \citep{schulman2017proximal}, where the reward model is trained from preference annotations, and an oracle variant, \textbf{PPO-GT}, which utilizes the actual reward function in the sentiment control scenario. Our sentiment experiments explore two PPO-GT implementations: a standard off-the-shelf solution \cite{leandro_von_werra_2023_7790115}, and an enhanced version with reward normalization and further hyperparameter optimization to boost performance. These modifications are also incorporated in PPO with learned rewards. Finally, we implement the \textbf{Best of $N$} strategy as a strong baseline. Here, $N$ candidate responses are generated using the SFT model (or Preferred-FT in the dialogue setup), and the candidate with the highest predicted reward—according to the preference-trained reward function—is selected. While this approach provides a robust upper bound, its practical utility is limited by the significant computational resources required to generate and score $N$ completions per query during inference.

\subsection{How well can DPO optimize the RLHF objective?}

RLHF algorithms typically employ a KL-regularized reward maximization objective, which seeks to enhance reward acquisition while simultaneously penalizing deviations from a reference policy. Thus, comparisons across algorithms should jointly consider both the obtained reward and the associated KL divergence; a marginal increase in reward at the expense of a substantially higher KL may not constitute a genuine improvement. In Figure~\ref{fig:frontier-tldr-main}, we display the trade-off between reward and KL (the reward-KL frontier) across different methods in a sentiment classification scenario. For each algorithm, we perform several training experiments, varying the hyperparameters governing policy conservativeness (target KL $\in\{3,6,9,12\}$ for PPO, $\beta \in \{0.05,0.1,1,5\}$ for DPO, $\alpha\in\{0.05,0.1,0.5,1\}$ for unlikelihood, and distinct random seeds for preferred-FT), resulting in a total of 22 distinct runs. After every increment of 100 training steps, and continuing until convergence, we assess each trained policy on a suite of test prompts, calculating both the mean reward under the actual reward function and the average sequence-level KL divergence\footnote{Specifically, the sum over timesteps of per-step KL-divergences.} with respect to the reference policy, $\text{KL}\left(\pi\mid \mid \pi_\text{ref}\right)$. The findings indicate that DPO generates the most favorable reward-KL frontier, attaining high rewards without incurring substantial increases in KL. This outcome is particularly striking for several reasons. Notably, although DPO and PPO target the same objective, DPO achieves superior efficiency, strictly outperforming PPO along the reward/KL spectrum. Moreover, DPO continues to deliver a superior frontier compared to PPO, \emph{even under conditions where PPO has access to the true reward signal} (PPO-GT).

\subsection{Can DPO scale to real preference datasets?}
\label{sec:dpo-real-datasets}
We proceed to assess DPO's fine-tuning capabilities on tasks involving summarization and single-turn dialogue. In the case of summarization, traditional automatic evaluation criteria such as ROUGE are often only weakly aligned with human judgment~\citep{stiennon2022learning}, and previous research has demonstrated that leveraging PPO for LM fine-tuning based on human feedback can yield more satisfactory summaries. To compare different approaches, we generate completions on the test partition of the TL;DR summarization dataset and calculate their average win rates against reference completions. Sampling for each method is performed across a range of temperatures from 0.0 to 1.0; resulting win rates are depicted in Figure~\ref{fig:frontier-tldr-main} (right). All methods—DPO, PPO, and Preferred-FT—begin fine-tuning from the same GPT-J SFT model\footnote{\url{https://huggingface.co/CarperAI/openai_summarize_tldr_sft}}. Our results indicate that DPO attains a win rate near 61\% when sampled at temperature 0.0, outperforming PPO, which achieves around 57\% at its optimal sampling temperature of 0.0. Furthermore, DPO secures a greater peak win rate in comparison to the best-of-$N$ baseline. It should be emphasized that we did not perform extensive tuning of the $\beta$ hyperparameter for DPO, suggesting that these outcomes may not reflect its full capability. Additionally, DPO exhibits increased resilience to changes in sampling temperature relative to PPO, which sees performance deteriorate toward that of the baseline GPT-J at elevated temperatures. Preferred-FT yields only marginal gains over the original SFT model. Lastly, a direct human evaluation between DPO and PPO, presented in Section~\ref{sec:human-judgments}, shows that DPO completions sampled at temperature 0.25 were favored 58\% of the time over PPO completions generated at temperature 0.

For the single-turn dialogue task, we assess the various approaches using a subset of the test partition from the Anthropic HH dataset \citep{bai2022training}, specifically focusing on exchanges involving just one human-assistant turn. In the GPT-4 assessments, we determine each method’s win rate by referencing the preferred completions on the test set. Since there is no established SFT baseline for this task, we employ a pre-trained Pythia-2.8B model, fine-tune a reference model with Preferred-FT on the selected completions to ensure outputs are in-distribution, and subsequently apply DPO for further training. Additionally, we benchmark against the top result out of 128 Preferred-FT completions (having observed that increasing $N$ beyond 128 offers no further improvement for this setup; refer to Appendix Figure~\ref{fig:best-of-n}) as well as a 2-shot prompt variant using the base Pythia-2.8B model. Our findings indicate that DPO either matches or surpasses the best-performing temperature settings of other methods. 

We further evaluate an RLHF agent that was trained via PPO on the Anthropic HH dataset\footnote{\url{https://huggingface.co/reciprocate/ppo_hh_pythia-6B}}, sourced from a reputable implementation\footnote{\url{https://github.com/CarperAI/trlx/tree/main/examples/hh}}, but are unable to identify a prompt or sampling temperature configuration that outperforms the unadapted Pythia-2.8B model. Based on results from TL;DR and the shared reward objective across methods, we interpret the Best of 128 approach as a rough estimate of PPO-level outcomes. In summary, DPO stands out as the only method that is both computationally efficient and capable of outperforming the preferred completions in the Anthropic HH dataset, while achieving results on par with or superior to the computationally heavy Best of 128 baseline. Figure~\ref{fig:dialogue-main} further demonstrates that DPO attains optimal performance in a relatively short training time.

\subsection{Generalization to a new input distribution}

To investigate the robustness of PPO and DPO to distributional changes, we assess the performance of the policies trained in our Reddit TL;DR summarization study when applied to a new domain—news articles from the test partition of the CNN/DailyMail dataset \citep{nallapati-etal-2016-abstractive}. The evaluation uses the optimal sampling temperatures previously determined from TL;DR (0 and 0.25), with results summarized in Table~\ref{tab:ood}. We calculated the rate at which GPT-4 selected a model-generated summary over the reference summary, adapting the original Reddit TL;DR prompt (GPT-4 (C)) by substituting ``forum post'' with ``news article.'' DPO demonstrates superior performance relative to PPO in this shifted distribution context. These findings offer preliminary support for the conclusion that DPO-trained policies can achieve generalization on par with, or better than, those obtained with PPO—despite DPO not leveraging the extra unlabeled Reddit TL;DR prompts available to PPO.

\subsection{Validating GPT-4 judgments with human judgments}
\label{sec:human-judgments}
We undertake a human evaluation to test the consistency and credibility of GPT-4's judgments on summaries, focusing on outputs from the TL;DR experiment and employing two variants of the GPT-4 prompt. The \textbf{GPT-4 (S)} (simple) prompt requests only a judgment as to which summary better encapsulates the key information from the post. Conversely, the \textbf{GPT-4 (C)} (concise) prompt extends this by also considering the conciseness of the summary; this variant is used because, under the simple prompt, GPT-4 tends to favor summaries that are longer and more repetitive, which differs from human preferences. Full prompt texts are provided in Appendix~\ref{app:prompts}. Three systems are evaluated: the highest-performing (DPO, temperature 0.25), lowest-performing (PPO, temperature 1.0), and an intermediate (SFT, temperature 0.25), each compared against the best temperature PPO baseline (greedy sampling), in order to capture a range of output qualities. Our analysis shows that, for both prompts, GPT-4's preferences are aligned with those of human raters to a similar degree as human raters agree amongst themselves (for DPO and PPO-1 comparisons, where multiple human assessments were feasible). Notably, the \textbf{GPT-4 (C)} prompt generates win rates that better reflect human judgments, motivating its selection for reporting the core outcomes in Section~\ref{sec:dpo-real-datasets}. Appendix~\ref{app:human-study} provides further information regarding the human evaluation setup, the web interface, and a list of participating raters.

\section{Discussion}
Preference-based learning represents an effective and scalable approach for aligning and enhancing the capabilities of language models. In this work, we presented DPO, a straightforward method for training language models from preference data, circumventing the need for reinforcement learning. Instead of reframing the preference learning task to suit conventional RL frameworks and employing standard RL solutions, DPO establishes a direct correspondence between policy distributions of language models and reward formulations, thereby facilitating alignment with human preferences through a standard cross-entropy objective—entirely omitting reinforcement learning and without sacrificing expressiveness. Remarkably, with little to no hyperparameter adjustment, DPO attains comparable or superior results relative to established RLHF approaches, such as those leveraging PPO, which makes preference-based training significantly more accessible.

\textbf{Limitations \& Future Work.} Several open questions emerge from our study. One area for further exploration is the out-of-distribution generalization properties of policies learned by DPO compared to those optimized via explicit reward functions; preliminary findings imply similar generalization to PPO-driven methods, yet this warrants deeper investigation. For instance, could methods such as self-labeling using the DPO policy unlock effective utilization of unlabeled prompts? Another pertinent direction involves examining how direct preference optimization may contribute to phenomena like reward over-optimization, as hinted by the moderate dip in performance observed in Figure~\ref{fig:dialogue-main}-right—clarifying whether this is symptomatic of such effects. Additionally, although we have tested models up to the 6B parameter scale, extending DPO’s application to models many times larger remains a promising avenue for future research. As for evaluation procedures, we observed that GPT-4’s win rate assessments are sensitive to prompt formulation; subsequent studies could investigate optimal methods for eliciting reliable automated judgments. Finally, DPO’s utility may extend well beyond language modeling, offering potential in training generative systems for various data modalities.